% META: {"id": "1SPE-SECDEG-CO-001", "chapitre": "1SPE-SECOND-DEGRE", "type_objet": "corrige", "exercice_ref": "1SPE-SECDEG-EX-001", "capacites_codes": ["C1"], "sources_inspiration": [], "status": "generated"}
% BEGIN-VERIFY
% from sympy import symbols, expand, factor, simplify, Rational
% x = symbols('x')
% f = 2*x**2 - 3*x - 5
% # 1. Coefficients
% assert f.coeff(x, 2) == 2   # a = 2
% assert f.coeff(x, 1) == -3  # b = -3
% assert f.coeff(x, 0) == -5  # c = -5
% # 2. Sommet
% xs = Rational(-(-3), 2*2)   # -b/(2a) = 3/4
% assert xs == Rational(3, 4)
% beta = f.subs(x, xs)
% assert beta == Rational(-49, 8)
% # 3. Forme canonique
% assert simplify(f - (2*(x - Rational(3,4))**2 - Rational(49,8))) == 0
% # 4. Forme factorisee
% assert simplify(f - 2*(x + 1)*(x - Rational(5,2))) == 0
% from sympy import expand as ex
% assert ex(2*(x + 1)*(x - Rational(5,2))) == 2*x**2 - 3*x - 5
% END-VERIFY
\begin{corrige}{1SPE-SECDEG-EX-001}

\textbf{1. Identification des coefficients.}

\commentaireMarge{Lire $a$, $b$, $c$ dans $ax^2+bx+c$.}

Le trinôme $f(x) = 2x^2 - 3x - 5$ est de la forme $ax^2 + bx + c$ avec :
\[a = 2, \quad b = -3, \quad c = -5.\]

\medskip
\textbf{2. Coordonnées du sommet.}

\commentaireMarge{Formule du sommet : $\alpha = -\dfrac{b}{2a}$.}

L'abscisse du sommet est :
\[\alpha = -\frac{b}{2a} = -\frac{-3}{2 \times 2} = \frac{3}{4}.\]

L'ordonnée du sommet est :
\[\beta = f\!\left(\frac{3}{4}\right) = 2\times\left(\frac{3}{4}\right)^2 - 3\times\frac{3}{4} - 5 = 2\times\frac{9}{16} - \frac{9}{4} - 5 = \frac{9}{8} - \frac{18}{8} - \frac{40}{8} = -\frac{49}{8}.\]

Le sommet est $S\!\left(\dfrac{3}{4}\,;\,-\dfrac{49}{8}\right)$.

\medskip
\textbf{3. Forme canonique.}

\commentaireMarge{$f(x)=a(x-\alpha)^2+\beta$.}

On a $a = 2$, $\alpha = \dfrac{3}{4}$, $\beta = -\dfrac{49}{8}$, donc :
\[f(x) = 2\left(x - \frac{3}{4}\right)^2 - \frac{49}{8}.\]

\medskip
\textbf{4. Vérification de la forme factorisée.}

\commentaireMarge{Développer et simplifier pour vérifier.}

On développe $2\!\left(x+1\right)\!\left(x - \dfrac{5}{2}\right)$ :
\[2\!\left(x+1\right)\!\left(x - \frac{5}{2}\right) = 2\!\left(x^2 - \frac{5}{2}x + x - \frac{5}{2}\right) = 2\!\left(x^2 - \frac{3}{2}x - \frac{5}{2}\right) = 2x^2 - 3x - 5.\]

On retrouve bien $f(x)$. \hfill $\square$

\end{corrige}
